\section{Performance dimension: Inter-node}

% Inter-node behaviour means codes that run with MPI, i.e.~employ multiple ranks.
% Before we start, we need a setup where we can alter the problem size easily (cmp.~Section~\ref{section:performance-assessment:preparation}) and understand the arising complexity exactly.
% From hereon, we refer exclusively to a normalised runtime $t^{\text{norm}}(p)_{s}$ which is the time that we need for a scenario $s$ to compute one subresult.
% This could be an update of a cell or a particle or any other reasonable quantity.
% 
% 
% I personally prefer to argue over inter-node efficiency through weak scaling.
% Our goal is to assess to which degree we could employ twice as many compute nodes and have twice the throughput.
% Many books refer to ``twice the problem size'', but this is misleading as twice as big problems might, on a fixed number of ranks, needs more than twice the time.
% So we stick to this normalised time from hereon.
% 
% 
% \begin{assessmentcrime}
%  A sequence $s_0, s_1, s_2, \ldots$ of scenarios with increasing problem sizes is used, but all of them would fit on one or few nodes. 
% \end{assessmentcrime}
% 
% 
% \noindent
% Indeed, the normalised timings for the problems however should stem from setups which really need the number of nodes.
% If you split up a problem which could easily fit on one single node into more and more subparts, you will only measure subpartitioning overhead. 
% If people suffer from such overhead, I would argue that this is a user error (although we use exactly this idea below to assess this propery).
% 
% 
% \subsection{Is there a problem?}
% 
% The underlying Gustafson's law~\cite{xxx} reads
% 
% \[
%  t(1) = p \cdot f \cdot t(p)  + (1-f) t(p),
% \]
% 
% \noindent
% but that's not that useful here.

